\documentclass[12pt]{article}
\usepackage{amsmath,amssymb}

\begin{document}
\section*{{2017 AIME II Problems/Problem 2}}
Source: AoPS Wiki

\subsection*{Solution}
There are two scenarios in which $T_4$ wins. The first scenario is where $T_4$ beats $T_1$ , $T_3$ beats $T_2$ , and $T_4$ beats $T_3$ , and the second scenario is where $T_4$ beats $T_1$ , $T_2$ beats $T_3$ , and $T_4$ beats $T_2$ . Consider the first scenario. The probability $T_4$ beats $T_1$ is $\frac{4}{4+1}$ , the probability $T_3$ beats $T_2$ is $\frac{3}{3+2}$ , and the probability $T_4$ beats $T_3$ is $\frac{4}{4+3}$ . Therefore the first scenario happens with probability $\frac{4}{4+1}\cdot\frac{3}{3+2}\cdot\frac{4}{4+3}$ . Consider the second scenario. The probability $T_4$ beats $T_1$ is $\frac{4}{1+4}$ , the probability $T_2$ beats $T_3$ is $\frac{2}{2+3}$ , and the probability $T_4$ beats $T_2$ is $\frac{4}{4+2}$ . Therefore the second scenario happens with probability $\frac{4}{1+4}\cdot\frac{2}{2+3}\cdot\frac{4}{4+2}$ . By summing these two probabilities, the probability that $T_4$ wins is $\frac{4}{4+1}\cdot\frac{3}{3+2}\cdot\frac{4}{4+3}+\frac{4}{1+4}\cdot\frac{2}{2+3}\cdot\frac{4}{4+2}$ . Because this expression is equal to $\frac{256}{525}$ , the answer is $256+525=\boxed{781}$ .

\end{document}
